\chapter{Dérivation}
\textit{Dériver, c'est mesurer la vitesse à laquelle une fonction varie au voisinage
d'un point. Cette idée — comparer une variation de sortie à une variation d'entrée
de plus en plus petite — débouche sur un outil d'une puissance remarquable : la
\emph{tangente} à une courbe et, surtout, un critère mécanique pour déterminer le
sens de variation d'une fonction et localiser ses extremums. Ce chapitre construit
le nombre dérivé, dresse le tableau des dérivées usuelles, donne les règles de
calcul, puis relie le signe de $f'$ aux variations de $f$.}

\section{Nombre dérivé en un point}

\begin{definition}[Taux d'accroissement]
Soit $f$ une fonction définie sur un intervalle $I$ et $a\in I$. Pour un réel $h\neq 0$
tel que $a+h\in I$, le \textbf{taux d'accroissement} de $f$ entre $a$ et $a+h$ est
\[ \tau(h)=\frac{f(a+h)-f(a)}{h}. \]
C'est la pente de la droite qui joint les points $A(a\,;\,f(a))$ et $M(a+h\,;\,f(a+h))$
de la courbe (droite \textbf{sécante}).
\end{definition}

\begin{definition}[Nombre dérivé]
On dit que $f$ est \textbf{dérivable en $a$} lorsque le taux d'accroissement admet une
\textbf{limite finie} quand $h$ tend vers $0$. Cette limite est appelée
\textbf{nombre dérivé} de $f$ en $a$ et notée $f'(a)$ :
\[ f'(a)=\lim_{h\to 0}\frac{f(a+h)-f(a)}{h}. \]
\end{definition}

\begin{remarque}
En posant $x=a+h$ (donc $h=x-a$ et $h\to 0 \iff x\to a$), on peut aussi écrire
\[ f'(a)=\lim_{x\to a}\frac{f(x)-f(a)}{x-a}. \]
Les deux formes sont équivalentes : on choisit la plus commode selon le calcul.
\end{remarque}

\begin{exemple}
Calculons le nombre dérivé de $f(x)=x^2$ en $a=3$. Pour $h\neq 0$ :
\[ \tau(h)=\frac{(3+h)^2-3^2}{h}=\frac{9+6h+h^2-9}{h}=\frac{6h+h^2}{h}=6+h. \]
Quand $h\to 0$, $\tau(h)\to 6$. Donc $f$ est dérivable en $3$ et $f'(3)=6$.
\end{exemple}

\section{Interprétation géométrique : la tangente}

\begin{aretenir}[De la sécante à la tangente]
Quand $h$ tend vers $0$, le point $M(a+h\,;\,f(a+h))$ se rapproche de $A(a\,;\,f(a))$
et la droite sécante $(AM)$ tend vers une position limite : la \textbf{tangente}
à la courbe $\mathcal{C}_f$ au point $A$. Le nombre dérivé $f'(a)$ est exactement
le \textbf{coefficient directeur} (la pente) de cette tangente.
\end{aretenir}

\begin{illus}
\begin{tikzpicture}[scale=1]
  \begin{axis}[
    width=10cm,height=7cm,
    axis lines=middle,
    xmin=-0.4,xmax=3.4,ymin=-0.6,ymax=5.4,
    xtick=\empty,ytick=\empty,
    clip=false,
    x label style={anchor=west}, y label style={anchor=south},
    xlabel={$x$},ylabel={$y$},
  ]
    \addplot[thick,onbo,domain=-0.3:2.3,samples=80]{x^2};
    % tangente en a=1 : y = 2(x-1)+1 = 2x-1
    \addplot[thick,onbgreen,domain=-0.2:2.6,samples=2]{2*x-1};
    % sécante entre a=1 et a+h=2 : pente 3, passe par (1,1) : y=3x-2
    \addplot[onbmuted,dashed,domain=0.2:2.8,samples=2]{3*x-2};
    \node[onbo,circle,fill,inner sep=1.4pt,label={below right:$A(a;f(a))$}] at (axis cs:1,1){};
    \node[onbmuted,circle,fill,inner sep=1.4pt,label={right:$M(a+h;f(a+h))$}] at (axis cs:2,4){};
    \node[onbgreen,anchor=west] at (axis cs:2.6,4.2){tangente};
  \end{axis}
\end{tikzpicture}
\fig{Figure — la sécante $(AM)$ (pointillés) pivote vers la tangente (verte) lorsque $h\to 0$ ; sa pente tend vers $f'(a)$.}
\end{illus}

\begin{theoreme}[Équation de la tangente]
Si $f$ est dérivable en $a$, la courbe $\mathcal{C}_f$ admet au point d'abscisse $a$
une tangente $\mathcal{T}$ d'équation
\[ \boxed{\,y=f'(a)\,(x-a)+f(a)\,.} \]
\end{theoreme}

\begin{remarque}
Cette droite passe par $A(a\,;\,f(a))$ (remplacer $x$ par $a$ donne $y=f(a)$) et a
pour pente $f'(a)$ : c'est bien la droite affine qui « épouse » la courbe en $A$.
Si $f'(a)=0$, la tangente est \textbf{horizontale}.
\end{remarque}

\begin{exemple}
Pour $f(x)=x^2$ au point d'abscisse $a=1$, on a $f(1)=1$ et $f'(1)=2$ (même calcul
qu'en $a=3$ : $\tau(h)=2+h\to 2$). La tangente a pour équation
\[ y=2(x-1)+1=2x-1. \]
\end{exemple}

\section{Fonction dérivée et dérivées de référence}

\begin{definition}[Fonction dérivée]
Si $f$ est dérivable en tout réel $x$ d'un intervalle $I$, on définit sur $I$ la
\textbf{fonction dérivée} de $f$, notée $f'$, qui à chaque $x$ associe le nombre
dérivé $f'(x)$.
\end{definition}

\begin{propriete}[Dérivées des fonctions de référence]
Le tableau ci-dessous donne la dérivée des fonctions usuelles, avec le domaine sur
lequel le calcul est valable.
\begin{center}\renewcommand{\arraystretch}{1.5}
\begin{tabular}{@{}l l l@{}}
\toprule
\textbf{Fonction $f(x)$} & \textbf{Dérivée $f'(x)$} & \textbf{Valable sur} \\
\midrule
$k$ (constante)        & $0$                       & $\R$ \\
$x$                    & $1$                       & $\R$ \\
$x^2$                  & $2x$                      & $\R$ \\
$x^n$\ \ ($n\in\N^*$)  & $n\,x^{\,n-1}$            & $\R$ \\
$\dfrac{1}{x}$         & $-\dfrac{1}{x^2}$         & $\R^*$ \\
$\sqrt{x}$             & $\dfrac{1}{2\sqrt{x}}$    & $]0\,;\,+\infty[$ \\
\bottomrule
\end{tabular}
\end{center}
\end{propriete}

\begin{theoreme}[Dérivée de $x^2$ (démonstration)]
La fonction $f(x)=x^2$ est dérivable sur $\R$ et $f'(x)=2x$.
\end{theoreme}

\begin{exemple}[Preuve]
Pour $x\in\R$ et $h\neq 0$ :
\[ \frac{f(x+h)-f(x)}{h}=\frac{(x+h)^2-x^2}{h}=\frac{2xh+h^2}{h}=2x+h. \]
Quand $h\to 0$, cette quantité tend vers $2x$. Donc $f'(x)=2x$.
\end{exemple}

\begin{remarque}
La fonction $x\mapsto\sqrt{x}$ est définie en $0$ mais \emph{n'y est pas dérivable} :
le taux $\dfrac{\sqrt{h}-0}{h}=\dfrac{1}{\sqrt{h}}$ tend vers $+\infty$ quand $h\to 0^+$.
Graphiquement, la courbe possède en $O$ une \textbf{tangente verticale}.
\end{remarque}

\section{Opérations sur les dérivées}

\begin{propriete}[Règles de dérivation]
Soit $u$ et $v$ deux fonctions dérivables sur un même intervalle $I$, et $k$ une
constante réelle. Sur $I$ (et là où le quotient est défini) :
\begin{center}\renewcommand{\arraystretch}{1.6}
\begin{tabular}{@{}l l@{}}
\toprule
\textbf{Fonction} & \textbf{Dérivée} \\
\midrule
$u+v$            & $u'+v'$ \\
$k\,u$           & $k\,u'$ \\
$u\,v$           & $u'v+u\,v'$ \\
$\dfrac{1}{v}$ \ ($v\neq 0$)   & $-\dfrac{v'}{v^2}$ \\
$\dfrac{u}{v}$ \ ($v\neq 0$)   & $\dfrac{u'v-u\,v'}{v^2}$ \\
$u^n$\ \ ($n\in\N^*$)          & $n\,u'\,u^{\,n-1}$ \\
\bottomrule
\end{tabular}
\end{center}
\end{propriete}

\begin{theoreme}[Dérivée d'un produit (démonstration)]
Si $u$ et $v$ sont dérivables, alors $uv$ l'est et $(uv)'=u'v+uv'$.
\end{theoreme}

\begin{exemple}[Idée de la preuve]
Formons le taux d'accroissement de $w=uv$ en $a$, puis ajoutons et retranchons un
terme intermédiaire :
\[ \frac{u(a+h)v(a+h)-u(a)v(a)}{h}
 =\frac{u(a+h)-u(a)}{h}\,v(a+h)+u(a)\,\frac{v(a+h)-v(a)}{h}. \]
Quand $h\to 0$ : le premier taux tend vers $u'(a)$, le second vers $v'(a)$, et
$v(a+h)\to v(a)$ (car $v$, dérivable, est continue). On obtient
$w'(a)=u'(a)v(a)+u(a)v'(a)$.
\end{exemple}

\begin{exemple}
Dérivons $f(x)=(2x^2-1)(x+3)$. Avec $u=2x^2-1$ ($u'=4x$) et $v=x+3$ ($v'=1$) :
\[ f'(x)=u'v+uv'=4x(x+3)+(2x^2-1)\times 1=4x^2+12x+2x^2-1=6x^2+12x-1. \]
\end{exemple}

\begin{exemple}
Dérivons $g(x)=\dfrac{x-1}{x^2+1}$. Avec $u=x-1$ ($u'=1$) et $v=x^2+1$ ($v'=2x$) :
\[ g'(x)=\frac{u'v-uv'}{v^2}=\frac{1\cdot(x^2+1)-(x-1)(2x)}{(x^2+1)^2}
 =\frac{-x^2+2x+1}{(x^2+1)^2}. \]
\end{exemple}

\section{Dérivée et sens de variation}

\begin{theoreme}[Signe de $f'$ et variations de $f$]
Soit $f$ une fonction dérivable sur un intervalle $I$.
\begin{itemize}
  \item Si $f'(x)>0$ pour tout $x$ de $I$, alors $f$ est \textbf{strictement croissante} sur $I$ ;
  \item si $f'(x)<0$ pour tout $x$ de $I$, alors $f$ est \textbf{strictement décroissante} sur $I$ ;
  \item si $f'(x)=0$ pour tout $x$ de $I$, alors $f$ est \textbf{constante} sur $I$.
\end{itemize}
\end{theoreme}

\begin{aretenir}[Étudier les variations d'une fonction]
On suit toujours le même plan : \textbf{1.} calculer $f'(x)$ ; \textbf{2.} étudier le
\emph{signe} de $f'(x)$ ; \textbf{3.} en déduire les variations dans un
\textbf{tableau de variations}. Le sens des flèches recopie le signe de $f'$.
\end{aretenir}

\begin{exemple}
Étudions $f(x)=x^2-2x+3$ sur $\R$. On a $f'(x)=2x-2=2(x-1)$, qui s'annule en $x=1$,
est négatif pour $x<1$ et positif pour $x>1$. Donc $f$ décroît sur $]-\infty\,;\,1]$
puis croît sur $[1\,;\,+\infty[$ ; $f(1)=2$.
\end{exemple}

\begin{illus}
\begin{tikzpicture}[scale=1,x=1.6cm,y=1cm]
  % cadre du tableau
  \draw[onbline] (0,0) rectangle (4,3);
  \draw[onbline] (0,2) -- (4,2);   % séparation x / f'
  \draw[onbline] (0,1) -- (4,1);   % séparation f' / f
  \draw[onbline] (1,3) -- (1,0);   % colonne d'étiquettes
  % étiquettes de lignes
  \node[onbink] at (0.5,2.5){$x$};
  \node[onbink] at (0.5,1.5){$f'(x)$};
  \node[onbink] at (0.5,0.5){$f$};
  % ligne des x
  \node[onbink] at (1.4,2.5){$-\infty$};
  \node[onbink] at (2.5,2.5){$1$};
  \node[onbink] at (3.7,2.5){$+\infty$};
  % ligne du signe de f'
  \node[onbmuted] at (1.8,1.5){$-$};
  \node[onbink]   at (2.5,1.5){$0$};
  \node[onbmuted] at (3.2,1.5){$+$};
  % variations
  \draw[onbmuted,->] (1.3,0.85) -- (2.4,0.2);
  \draw[onbmuted,->] (2.6,0.2) -- (3.7,0.85);
  \node[onbo] at (2.5,0.05){$2$};
\end{tikzpicture}
\fig{Figure — tableau de variations de $f(x)=x^2-2x+3$ : minimum $2$ atteint en $x=1$.}
\end{illus}

\section{Extremums}

\begin{definition}[Extremum local]
La fonction $f$ admet en $x_0$ un \textbf{maximum local} (resp. \textbf{minimum
local}) s'il existe un intervalle ouvert contenant $x_0$ sur lequel $f(x)\le f(x_0)$
(resp. $f(x)\ge f(x_0)$).
\end{definition}

\begin{theoreme}[Extremum et dérivée]
Soit $f$ dérivable sur un intervalle ouvert $I$ et $x_0\in I$. Si $f$ admet un
extremum local en $x_0$, alors $f'(x_0)=0$.

\medskip
\textbf{Réciproque pratique.} Si $f'$ \textbf{s'annule en $x_0$ en changeant de
signe}, alors $f$ présente un extremum en $x_0$ : un \textbf{maximum} si $f'$ passe
du $+$ au $-$, un \textbf{minimum} si $f'$ passe du $-$ au $+$.
\end{theoreme}

\begin{remarque}
L'annulation de $f'$ \emph{seule} ne suffit pas : pour $f(x)=x^3$, $f'(0)=0$ mais $f'$
reste positive de part et d'autre, donc $f$ n'a \textbf{pas} d'extremum en $0$
(c'est un point d'inflexion à tangente horizontale). Il faut bien un \emph{changement
de signe}.
\end{remarque}

\begin{exemple}
Pour $f(x)=x^2-2x+3$ (ci-dessus), $f'$ passe du $-$ au $+$ en $x=1$ : $f$ admet donc
un \textbf{minimum} égal à $f(1)=2$, atteint en $x=1$.
\end{exemple}

% ════════════════════════════════════════════════════════════
\section{L'essentiel à retenir}

\begin{synthese}
\textbf{Nombre dérivé.} $f'(a)=\displaystyle\lim_{h\to 0}\dfrac{f(a+h)-f(a)}{h}$ :
c'est la pente de la tangente à $\mathcal{C}_f$ au point d'abscisse $a$.

\smallskip
\textbf{Tangente.} $\ y=f'(a)(x-a)+f(a)$. Horizontale lorsque $f'(a)=0$.

\smallskip
\textbf{Dérivées usuelles.}
\begin{center}\renewcommand{\arraystretch}{1.35}
\begin{tabular}{@{}l l l l@{}}
\toprule
$f(x)$ & $f'(x)$ & $f(x)$ & $f'(x)$ \\
\midrule
$k$    & $0$         & $\dfrac{1}{x}$  & $-\dfrac{1}{x^2}$ \\
$x^n$  & $n\,x^{n-1}$ & $\sqrt{x}$     & $\dfrac{1}{2\sqrt{x}}$ \\
\bottomrule
\end{tabular}
\end{center}

\smallskip
\textbf{Opérations.}
\[ (u+v)'=u'+v';\quad (ku)'=ku';\quad (uv)'=u'v+uv'; \]
\[ \Big(\tfrac1v\Big)'=-\tfrac{v'}{v^2};\quad \Big(\tfrac{u}{v}\Big)'=\tfrac{u'v-uv'}{v^2};
   \quad (u^n)'=n\,u'\,u^{n-1}. \]

\smallskip
\textbf{Variations.} $f'>0\Rightarrow f\nearrow$ ; $f'<0\Rightarrow f\searrow$ ;
$f'=0$ sur $I\Rightarrow f$ constante.

\smallskip
\textbf{Extremum.} en un extremum intérieur $f'$ s'annule ; réciproquement, un
\emph{changement de signe} de $f'$ donne un extremum ($+\!\to\!-$ : maximum,
$-\!\to\!+$ : minimum).
\end{synthese}

% ════════════════════════════════════════════════════════════
\section{Méthodes \& savoir-faire}

\methode{Calculer un nombre dérivé par la définition}
Former $\tau(h)=\dfrac{f(a+h)-f(a)}{h}$, développer et \emph{simplifier} par $h$,
puis faire tendre $h$ vers $0$.
\begin{exemple}
$f(x)=\dfrac1x$ en $a=2$ : $\tau(h)=\dfrac{\frac1{2+h}-\frac12}{h}
=\dfrac{2-(2+h)}{2(2+h)h}=\dfrac{-1}{2(2+h)}\xrightarrow[h\to0]{}-\dfrac14$. Donc
$f'(2)=-\dfrac14$ (cohérent avec $f'(x)=-\frac1{x^2}$).
\end{exemple}

\methode{Dériver une fonction polynôme ou rationnelle}
Identifier les morceaux, appliquer somme, produit par une constante, et la formule
$(x^n)'=n x^{n-1}$ ; pour un quotient, poser $u$ et $v$ puis appliquer
$\big(\frac uv\big)'=\frac{u'v-uv'}{v^2}$.
\begin{exemple}
$f(x)=3x^4-5x^2+x-7$ donne $f'(x)=12x^3-10x+1$.
\end{exemple}

\methode{Dériver une puissance $u^n$}
Reconnaître la forme $u^n$, calculer $u'$, puis appliquer $(u^n)'=n\,u'\,u^{n-1}$
sans développer.
\begin{exemple}
$f(x)=(3x-1)^4$ : $u=3x-1$, $u'=3$, donc $f'(x)=4\times 3\times(3x-1)^3=12(3x-1)^3$.
\end{exemple}

\methode{Déterminer l'équation d'une tangente}
Calculer $f(a)$ et $f'(a)$, puis remplacer dans $y=f'(a)(x-a)+f(a)$.
\begin{exemple}
$f(x)=x^3$ en $a=2$ : $f(2)=8$, $f'(x)=3x^2$ donc $f'(2)=12$. Tangente :
$y=12(x-2)+8=12x-16$.
\end{exemple}

\methode{Dresser un tableau de variations}
Calculer $f'(x)$, résoudre $f'(x)=0$, étudier le \emph{signe} de $f'$ entre ses
racines (signe d'un produit/quotient), puis reporter flèches et valeurs extrêmes.
\begin{exemple}
$f(x)=x^3-3x$ : $f'(x)=3x^2-3=3(x-1)(x+1)$, positif hors de $[-1\,;\,1]$, négatif
dedans. Donc $f$ croît sur $]-\infty\,;\,-1]$, décroît sur $[-1\,;\,1]$, croît sur
$[1\,;\,+\infty[$.
\end{exemple}

\methode{Localiser et justifier un extremum}
Au point où $f'$ s'annule, vérifier qu'elle \emph{change de signe} ; conclure
maximum ($+\to-$) ou minimum ($-\to+$), et donner sa valeur $f(x_0)$.
\begin{exemple}
Pour $f(x)=x^3-3x$, $f'$ passe du $+$ au $-$ en $x=-1$ : \textbf{maximum local}
$f(-1)=2$ ; du $-$ au $+$ en $x=1$ : \textbf{minimum local} $f(1)=-2$.
\end{exemple}
